\documentclass[a4paper,10pt]{llncs}
\usepackage[utf8]{inputenc}
\usepackage[T1]{fontenc}%
\usepackage{amssymb}
\usepackage{array}
\usepackage{amsmath}
\usepackage{graphicx}
\usepackage{cite}
\usepackage{color}
\usepackage{makecell}
%\usepackage{amsthm}
%opening
\title{}
\author{}

% \newtheorem{definition}{Definition}
% \newtheorem{proposition}{Proposition}
% \newtheorem{corollary}{Corollary}
% \newtheorem{lemma}{Lemma}
% \newtheorem{example}{Example}
% \newtheorem{theorem}{Theorem}

\def\E {{\mathbb E}}
\def\EE {{\mathbb E}}
\def\FF {{\mathbb F}}
\def\II {{\mathbb I}}
\def\LL {{\mathbb L}}
\def\NN {{\mathbb N}}
\def\PP {{\mathbb P}}
\def\RR {{\mathbb R}}
\def\ZZ {{\mathbb Z}}
\def\RRpos {{\mathbb R_{\geq 0}}}
\def\RRposi {{\mathbb R_{\geq 0}^{\infty}}}

\definecolor{purple}{rgb}{0.8, 0, 0.8}
\newcommand\todo[1]{{\color{red}\textbf{[TO DO:  #1]}}}
\newcommand\idea[1]{{\color{purple}\textbf{[IDEA:  #1]}}}
\begin{document}


%\maketitle
\begin{titlepage}
 \begin{center}
  % Chercher des logos
  \includegraphics[scale=0.17]{ENS.jpg}
  \includegraphics[scale=0.17]{Rennes.png}
  \includegraphics[scale=0.4]{rwth.jpg}\\[2cm]
  %\includegraphics[scale=0.5]{rwth.jpg}\\[2cm]

  \textsc{\LARGE ENS Rennes - University of Rennes 1}\\[1cm]
  \textsc{\Large Internship report: Master Degree, first year}\\[3.5cm]
  %\textsc{\Large Parcours Recherche et Innovation}\\[3.5cm]
  
  
  \textnormal{\LARGE Possibility Distribution Semantics for Probabilistic Programs with Nondeterminism}\\[2cm]
  
  \textnormal{Intern: Joshua Peignier}\\
  \textnormal{Supervisors: Benjamin Kaminski, Christoph Matheja}\\
  \textnormal{Team: Software Modeling and Verification}\\
  \textnormal{Institute: RWTH Aachen University (Aachen, Germany)}\\
  \textnormal{Dates: 15/05/2017 - 11/08/2017}\\[3cm]
 \end{center}
 

\end{titlepage}

 \begin{abstract}
In contrast to ordinary programs, probabilistic programs compute a probability distribution of output states for each given input state. One benefit of adding randomness into programs is that computationally hard problems, such as matrix multiplication or leader election protocols, can be solved (on average) more efficiently. A frequently used design concept to model unknown or overly complex probability distributions is nondeterministic choice. However, having probabilistic and nondeterministic choice within the same program (or even within loops) leads to subtle semantical intricacies. The goal of this report is to capture the semantics of both concepts uniformly using possibility theory. This would allow to simplify existing weakest-precondition style calculi for reasoning about probabilistic and nondeterministic programs.
  \begin{description}
  \item[Keywords:]Fuzzy imperative languages ; Fuzzy relations ; Possibility theory ; Weakest-precondition
  \end{description}
\end{abstract}

\section{Introduction}
\label{sec:intro}
Probabilistic programs are a particular class of programs, in which probabilistic choices are involved. Probabilistic choice consist of an instruction of the form $\{S_1\} [p] \{S_2\}$, where a program coming to this instruction has a probability $p$ of executing $S_1$ and a probability $1-p$ of executing $S_2$. For example, the simple program $\{ x := 2 \} [\frac{1}{3}] \{ x := 5 \}$ is a program that has a probability $\frac{1}{3}$ of setting the program variable $x$ to $2$, and a probability $\frac{2}{3}$ of setting $x$ to $5$. What makes this kind of program different from classical deterministic programs is that executing such a program with always the same input will not always result in the same output. Though they seem more complex than deterministic programs, probabilistic programs are useful in many ways. For example, it is simpler to solve some problems with probabilistic programs than with deterministic ones, as we can build probabilistic programs which have a better average complexity than deterministic programs (sometimes at the cost of a small probability of error). Classical examples of such programs are the quicksort algorithm, or primality tests, such as the Miller-Rabin test and the Solovay-Strassen test (for positive numbers), or the Berlekamp test (for polynomials). There also exists probabilistic algorithms to solve the leader election problem.\newline 
Besides probabilistic choice, there exists a frequently used concept in nondeterministic programs, called nondeterministic choice. One example is the following program: $\{ x := 2 \} [\!] \{ x := 5 \}$. In this example, the program will set $x$ nondeterministically to $2$, or to $5$, but we cannot assign a probability to any of these choices. More precisely, it makes no sense to assign them a probability. The nondeterministic choice is used in some programs where the outcomes of a choice are known, but making the mecanism behind the choices is either unknown, or uses overly complex probability distributions. \todo{Find a source} This instruction is used, for instance, in the model-checker Spin, to model the unpredictable behavior of a program where several choices are possible.\newline
\idea{Merge sections 1 and 2. And here, explain Hoare triples, Dijkstra's view of predicate transformers, etc...}
However, several major problems arise with the use of nondeterministic choice, the first one being the definition of proper semantics for this instruction. For instance, when trying to define predicate transformer semantics (i.e. when trying to define a program $P$ as a predicate transformer $T$ in the sense of Dijkstra, in \cite{Dijkstra76}, which transforms a given postcondition $\varphi$ into the corresponding weakest-postcondition $T(\varphi)$ such that the Hoare triple $\{T(\varphi)\}P\{\varphi\}$ is valid), several possible visions of the nondeterministic choices exist. The authors in \cite{McIver05,WuChen11} defined a weakest-precondition calculus on languages inspired from the pGCL language (an extended version of the Guarded Command Language from Dijkstra, including both type of choices mentioned above), where they consider that, for a program of the form $S_1 [\!] S_2$ and a postcondition $E$, the corresponding weakest precondition is considered as a minimum (or intersection) of the weakest preconditions of $E$ in $S_1$ and $S_2$, contrary to the work presented in \cite{WuChen08,WuChen12}, where the authors considered the same problem and defined the weakest precondition of $E$ in $S_1 [\!] S_2$ as the maximum (or union) of both.\newline
Another constraining problem is the complication arising from combination of probabilistic and nondeterministic choices in the same program.\todo{Give more details}\newline
The authors in \cite{WuChen08,WuChen11,WuChen12} attempted to define simpler semantics for fuzzy languages (detailed in \cite{Bueno93,Bueno97}) including the nondeterministic choice, using possibility theory (an alternative concept to probability theory to model uncertainty, detailed in \cite{Agarwal15,Shapiro09}). However, in none of their publications did they consider languages including probabilistic choice, as probabilities and possibilities are hard to combine. The authors in \cite{McIver05} tried to define a predicate transformer semantics for the pGCL language, including both type of choices, but got complicated semantics, where the weakest precondition we can deduce have harsh expressions.\newline
In this report, our goal is to give our own semantics of languages including both type of choices, and more precisely an expected value semantics (which can be derived from predicate transformer semantics), where we consider expected values of fuzzy random variables (first defined in \cite{PuriRal86}, and then simplified in \cite{Shapiro09}). It is a way to combine the work in \cite{WuChen08,WuChen11,WuChen12} using possibility theory with the languages presented in \cite{McIver05}.\newline
\todo{Insert a quick summary here, reminding what to find in each section}

\section{Context and Motivation}
\label{sec:context}
In this whole section, we will consider the pGCL language, presented in \cite{McIver05}, and defined by the following grammar\footnote{The original pGCL included an \texttt{abort} instruction corresponding to an error case. We chose, not to include it, as it can be represented by $\texttt{while} (\texttt{true})\{ \texttt{skip} \}$}:

\begin{align*}
 S ::= & \texttt{ skip } \,|\, x := A \,|\, S;S \,|\, S [p] S \,|\, S [\!] S \,|\ \\
 & \texttt{ if } (b) \texttt{ then } \{ S \} \texttt{ else } \{ S \} \,|\, \texttt{ while }(b) \texttt{ do }\{S\} \\
\end{align*}

This language is sufficient to describe any probabilistic or nondeterministic program. The \texttt{if\dots then \dots else} and \texttt{while} structure are the usual control structures encountered in many languages. The \texttt{skip} instruction corresponds to an empty instruction, where nothing is done (it is included in the language so that we do not have to create a specific \texttt{if\dots then \dots} instruction for the case where no \texttt{else} is needed.) $x := A$ is an assignment statement, where the program variable $x$ is set to the value of the expression $A$ in the current state. $S_1 ; S_2$ is a sequence assignment, modeling the execution of $S_1$ followed by $S_2$. We take interest primarly in the two remaining instructions.\newline
The $S_1 [p] S_2$ instruction models the fact that there is a probability $p$ that $S_1$ is executed, and a probability $1-p$ that $S_2$ is executed. Finally, the $S_1 [\!] S_2 $ (first introduced by Dijkstra in \cite{Dijkstra76}) is a statement that will nondeterministically execute either $S_1$ or $S_2$. The difference with the probabilistic choice is that no probability can be assigned to $S_1$ or $S_2$ because it makes no sense. This means that if $S_1 [\!] S_2 $ is executed a large number of times, no pattern will appear, whereas the execution of $S_1 [p] S_2$ a large number of times will tend to have a proportion $p$ of cases where $S_1$ is executed, and a proportion $1-p$ of cases where $S_2$ is executed.\newline
Note that we are using choice instructions with only two choices, but combination of instructions make possible the realisation of choices with any finite number (and even a countable number when using the loop) of outcomes.\newline

Contrary to the probabilistic choice, which is used in several algorithms in order to solve a given problem with a better average complexity than deterministic algorithms, the nondeterministic choice seems to be of little interest and looks like an unnecessary complication. But this operator managed to find its way in many publications which took interest in it (\cite{WuChen08,WuChen11,WuChen12,McIver05}), as it can be used in some situations, where a choice has to be made, where the outcomes of the choice are known, but where the mecanism hiding behind the choice is unknown (in this case, we are unable to assign probabilities), or where the mecanism uses an overly complex probability distribution.\todo{example of unknown or too complex mecanism}.\newline

Now, we consider the following pGCL program $P$: $\{ x := 2 [\!] x := 5 \} [p] \{ x := 7 \}$ (with $p \in ]0,1[$). \todo{apply the semantics of each paper to P to show the problem}

\section{Related Work}
\label{sec:related}

\section{Preliminaries}
\label{sec:preliminaries}

\section{Contribution}
\label{sec:contribution}

In this section, we consider programs over the language presented in Section \ref{sec:context}, and our goal is to define clearer semantics for these programs as the semantics we presented before. Recall that we showed in Section \ref{sec:preliminaries} that weakest-precondition calculus could be used in order to compute expected values of random variables. Our main idea here is to combine the concept of \textit{Fuzzy random variable} (introduced in \cite{PuriRal86}, and simplified in \cite{Shapiro09}) with the computation of expected values.\newline
Let $\Sigma$ be the set of states. We consider that the variables in our programs can only take \underline{real positive} values.\footnote{This is due to constraints explained further. But note that we can implement an arbitrary real variable by two variables, one for its absolute value, and one for its sign.} We adapted the following definition from \cite{PuriRal86}.

\begin{definition}{\textnormal{(Fuzzy Random Variable)\newline}}
 A Fuzzy Random Variable (short: FRV) is a mapping $X$ from the set of states $\Sigma$ to the powerset $P(\RRposi)$. This means that, if $X$ is an FRV, then for all $\sigma \in \Sigma$, $X(\sigma)$ is a subset of $\RRposi$. In fact, it is a set of possible values for $X$ in the state $\sigma$.
\end{definition}
\begin{example}
 In the following, for each program variable $x$, we denote by $\underline{x}$ the FRV such that, for all $\sigma \in \Sigma$, $\underline{x}(\sigma) = \{ x_{\sigma} \}$, meaning that $\underline{x}$ associates to each state $\sigma$ the singleton containing the value $x_\sigma$, which is the value of the program variable $x$ in the set $\sigma$. (Recall that each state is characterized uniquely by a tuple containing the values of each variable in a given order.). Therefore, we can say that the FRV $\underline{x}$ models the program variable $x$.
\end{example}

A fuzzy random variable $X$ over a program can model several concepts. For instance, it can be used to model the runtime of a program executed in a state $\sigma$ (which is why we chose to include $\infty$); in this example, $X(\sigma)$ would give all possible runtimes for the execution of one program starting from $\sigma$. It can also model the value of a program variable (or any function over the program variables) after the execution of a program.\newline
% For instance, for any program variable $x$, with the previous notations, for all $\sigma \in \Sigma$, $\underline{x}(\sigma) = \{ x \}$ is the expected value of the program variable $x$ after executing the program $\texttt{skip}$ in $\sigma$.\newline
% As all previous computations can be obtained from program variables (for instance, the runtime of a program (in the sense of "number of instructions executed before termination") can be obtained by adding a program variable $t$ representing the current runtime, initially $0$, and incrementing it at each instruction), we will in the following only consider computation of expected values of fuzzy random variables characterizing the value of a program variable after the execution of the program in one state.\newline

Let $F = \{f \,|\, f : \sigma \rightarrow P(\RRposi) \}$ be the set of fuzzy random variables, and $Progs$ be the set of pGCL programs. We define the following application:

\begin{definition}
 $ev : (Progs) \rightarrow (F \rightarrow F)$ is an application mapping each program $S$ to its FRV transformer semantics. This means that, for all $S \in Progs$, $ev[S]$ transforms a FRV $X \in F$ in another FRV $ev[S](X)$.
\end{definition}
Concretely, if $X$ is a FRV (i.e. an application mapping each state $\sigma$ to a subset $X(\sigma)$ of $\RRposi$), then $ev[S](X)(\sigma)$ is the set of possible expected values of the FRV $X$ after executing $S$ in the state $\sigma$.\newline
For instance, if we consider a program variable $x$, then the FRV $\underline{x}$ maps each state $\sigma$ to the set $\{x_\sigma\}$. Our goal is that the FRV $\underline{x}$ is transformed by $ev[S]$ into the FRV $ev[S](\underline{x})$, which maps each $\sigma$ to the set $ev[S](\underline{x})(\sigma)$ of possible values of the variable $x$ after executing $S$ in the state $\sigma$. Therefore, we can say that $ev[S](\underline{x})(\sigma)$ is the \emph{expected value} (or rather \emph{set of possible expected values}) of the FRV $\underline{x}$ after executing $S$ in the state $\sigma$; or, with other words, the set of possible expected values of the program variable $x$ after executing $S$ in $\sigma$.\newline


We give rules in Table \ref{table:rules_ev} of a predicate-transformer-style calculus for the computation of expected values.\newline
\begin{table}
\begin{center}
\begin{tabular}{|p{3cm}|p{9cm}|}
 \hline
 \thead{$S$} & \thead{$ev[S](X)$} \\
 \hline
 \thead{\texttt{skip}} & \thead{$X$} \\
 \hline
 \thead{$y := A$} & \thead{$\lambda\sigma.X(\sigma[y/A])$} \\
 \hline
 \thead{$S_1 ; S_2$} & \thead{$ev[S_1](ev[S_2](X))$} \\
 \hline
 \thead{$S_1 [p] S_2$} & \thead{$\lambda\sigma.\{t_1 p+t_2(1-p) \,|\, t_1 \in ev[S_1](X)(\sigma), t_2 \in ev[S_2](X)(\sigma) \}$} \\
 \hline
 \thead{$S_1 [\!] S_2$} & \thead{$\lambda\sigma. ev[S_1](X)(\sigma) \cup ev[S_2](X)(\sigma)$} \\
 \hline
 \thead{$\texttt{if } (b) \texttt{ then } \{ S_1 \}$ \\ $\texttt{ else } \{ S_2 \}$} & \thead{$\lambda\sigma.\{[\![b : true ]\!](\sigma)\cdot t_1 + [\![b : false ]\!](\sigma)\cdot t_2 \,|\,$ \\$t_1 \in ev[S_1](X)(\sigma), t_2 \in ev[S_2](X)(\sigma) \}$} \\
 \hline
 \thead{$\texttt{while }(b) \texttt{ do }\{S\}$} & \thead{lfp ($\lambda Y. (\lambda \sigma. \{[\![b : true ]\!](\sigma)\cdot t_1 + [\![b : false ]\!](\sigma)\cdot t_2 \,|\,$\\$t_1 \in ev[S](Y)(\sigma), t_2 \in X(\sigma) \}$))} \\
 \hline
\end{tabular}
\end{center}
\label{table:rules_ev}
\caption{Rules for defining the FRV transformer $ev$}
\end{table}

Consider the guideline example $P$: $\{ x := 2 [\!] x := 5 \} [p] \{ x := 7 \}$. We want to compute the expected value of $x$ after the execution of $P$. Thus, we have to determine $ev[P](\underline{x})$,which will map each state $\sigma$ to the expected value of $x$ after executing $P$ in $\sigma$. It is a function of $\sigma$, but the form of $P$ gives the intuition that it will be a constant function, not depending on $\sigma$.\newline
With the rules presented in Table \ref{table:rules_ev}, we can see that this requires first to compute the functions $ev[x := 2 [\!] x := 5](\underline{x})$ and $ev[x := 7](\underline{x})$; besides, the former requires to compute $ev[x := 2](\underline{x})$, $ev[x := 5](\underline{x})$.\newline

We get that:
$$ev[x := 2](\underline{x}) = \lambda\sigma.\underline{x}(\sigma[x/2]) = \lambda\sigma.\{x_{\sigma[x/2]}\} = \lambda\sigma.\{2\}$$ 
Indeed, the value of $x$ in $\sigma[x/2]$ is necessarily $2$. Recall that $\sigma[x/2]$ is the state obtained after setting $x$ to $2$ in the state $\sigma$. This means that the expected value of $x$ after executing $x := 2$ in any state $\sigma$ can only be $2$. By the same computation, we get $ev[x := 5](\underline{x}) = \lambda\sigma.\{5\}$ and $ev[x := 7](\underline{x}) = \lambda\sigma.\{7\}$.\newline

Now, we can compute that:
\begin{align*}
ev[x := 2 [\!] x := 5](\underline{x}) &= \lambda\sigma. ev[x := 2](\underline{x})(\sigma) \cup ev[x := 5](\underline{x})(\sigma) \\
&= \lambda\sigma. \{2\} \cup \{5\} = \lambda\sigma.\{2,5\} 
\end{align*}
This means that the possible expected values of $x$ after executing $x := 2 [\!] x := 5$ in any state $\sigma$ are $2$ and $5$.\newline

Finally, 
\begin{align*}
ev[P](\underline{x}) &= \lambda\sigma.\{t_1 p+t_2(1-p) \,|\, t_1 \in ev[x := 2 [\!] x := 5](\underline{x})(\sigma), t_2 \in ev[x := 7](\underline{x})(\sigma) \} \\
&= \lambda\sigma.\{t_1 p+t_2(1-p) \,|\, t_1 \in \{2,5\}, t_2 \in \{7\} \} \\
&= \lambda\sigma. \{2p+7(1-p),5p+7(1-p) \} 
\end{align*}



\section{Conclusion}
\todo{Explain that, for a program $S$, by considering each of its variables $x$ and computing $ev[S](\underline{x})$, we are able to fully characterize the program $S$ and get an understanding of what it does. Explain that it is clearer than what was done before.}
\label{sec:conclusion}


\bibliography{ref1}{}
\bibliographystyle{plain}

\end{document}
